% Part: sets-functions-relations
% Chapter: relations
% Section: special-properties

\documentclass[../../../include/open-logic-section]{subfiles}

\begin{document}

\olfileid{sfr}{rel}{prp}
\olsection{సంబంధాల ప్రత్యేక ధర్మాలు}

\begin{intro}
కొన్ని రకాల సంబంధాలు చాలా తరచుగా కనిపిస్తాయి కనుక వాటికి ప్రత్యేక
పేర్లు పెట్టారు. ఉదాహరణకు $\le$, $\subseteq$ తమతమ ప్రవేశాలలోని
మూలకాలను ఒకే విధంగా సంబంధపరుస్తాయి
($\le$కు $\Nat$నూ, $\subseteq$కు $\Pow{A}$నూ తీసుకోవచ్చు).
ఈ సంబంధాల మధ్య పోలికలు, తేడాలు కచ్చితంగా ఏమిటో తెలుసుకోవడానికి,
సంబంధాలకు ఉండగల కొన్ని ప్రత్యేక ధర్మాల ఆధారంగా వాటిని వర్గీకరిస్తాం.
వీటిలో కొన్ని ధర్మాలు, వాటి కలయికలు ప్రత్యేకించి ముఖ్యమైనవి:
అవి క్రమాలు, తుల్యతా సంబంధాలకు దారితీస్తాయి.
\end{intro}

\begin{defn}[స్వావర్తనత్వం]
ప్రతి $x \in A$కూ $Rxx$ కావడం అనేది $R \subseteq A^2$
అనే సంబంధం \emph{స్వావర్తన} సంబంధం కావడానికి అవసరమైన మరియు సరిపడే షరతు.
\end{defn}

\begin{defn}[సంక్రామకత్వం]
$Rxy$, $Ryz$ ఉన్నప్పుడల్లా $Rxz$ కూడా ఉండటం అనేది
$R \subseteq A^2$ అనే సంబంధం \emph{సంక్రామక} సంబంధం కావడానికి
అవసరమైన మరియు సరిపడే షరతు.
\end{defn}

\begin{defn}[సౌష్ఠవం]
$Rxy$ ఉన్నప్పుడల్లా $Ryx$ కూడా ఉండటం అనేది
$R \subseteq A^2$ అనే సంబంధం \emph{సౌష్ఠవ} సంబంధం కావడానికి
అవసరమైన మరియు సరిపడే షరతు.
\end{defn}

\begin{defn}[ప్రతిసౌష్ఠవం]
$Rxy$, $Ryx$ రెండూ ఉన్నప్పుడల్లా $x=y$ కావడం అనేది
$R \subseteq A^2$ అనే సంబంధం \emph{ప్రతిసౌష్ఠవ} సంబంధం కావడానికి
అవసరమైన మరియు సరిపడే షరతు (మరోలా చెప్పాలంటే:
$x\neq y$ అయితే $\lnot Rxy$ లేదా $\lnot Ryx$).
\end{defn}

\begin{explain}
సౌష్ఠవ సంబంధంలో $Rxy$, $Ryx$ ఎప్పుడూ రెండూ ఉంటాయి లేదా రెండూ ఉండవు.
ప్రతిసౌష్ఠవ సంబంధంలో $Rxy$, $Ryx$ రెండూ ఉండాలంటే $x = y$ కావాలి.
అయితే $x = y$ అయినప్పుడు $Rxy$, $Ryx$ తప్పక ఉండాలని ఇది
\emph{కోరదు}; ఆ అవకాశాన్ని నిరాకరించదంతే.
కనుక ప్రతిసౌష్ఠవ సంబంధం స్వావర్తన సంబంధం కావచ్చు, కానీ ప్రతి
ప్రతిసౌష్ఠవ సంబంధమూ స్వావర్తన సంబంధం కాదు.
అలాగే ప్రతిసౌష్ఠవంగా ఉండటం, కేవలం సౌష్ఠవంగా లేకపోవడం వేర్వేరు షరతులు.
నిజానికి ఒక సంబంధం ఒకేసారి సౌష్ఠవంగానూ ప్రతిసౌష్ఠవంగానూ ఉండవచ్చు
(తాదాత్మ్య సంబంధం అలాంటిదే).
\end{explain}

\begin{defn}[సంయుక్తత్వం]
ప్రతి $x,y\in A$కూ, $x \neq y$ అయితే $Rxy$ లేదా $Ryx$
ఉన్నప్పుడు $R \subseteq A^2$ అనే సంబంధాన్ని \emph{సంయుక్త}
సంబంధం అంటాం.
\end{defn}

\begin{prob}
ఈ ధర్మాలు గల సంబంధాలకు ఉదాహరణలు ఇవ్వండి:
(a) స్వావర్తన, సౌష్ఠవ ధర్మాలు ఉండి సంక్రామక ధర్మం లేకపోవడం;
(b) స్వావర్తన, ప్రతిసౌష్ఠవ ధర్మాలు ఉండటం;
(c) ప్రతిసౌష్ఠవ, సంక్రామక ధర్మాలు ఉండి స్వావర్తన ధర్మం లేకపోవడం;
(d) స్వావర్తన, సౌష్ఠవ, సంక్రామక ధర్మాలు ఉండటం.
సంఖ్యలపై లేదా సమితులపై సంబంధాలను ఉపయోగించవద్దు.
\end{prob} 
 
\begin{defn}[అస్వావర్తనత్వం]
ప్రతి $x \in A$కూ $Rxx$ లేకపోతే,
$R \subseteq A^2$ అనే సంబంధాన్ని \emph{అస్వావర్తన} సంబంధం అంటాం.
\end{defn}

\begin{defn}[ఏకదిశత్వం]
ఏ జత $x,y\in A$కూ $Rxy$, $Ryx$ రెండూ లేకపోతే,
$R \subseteq A^2$ అనే సంబంధాన్ని \emph{ఏకదిశ} సంబంధం అంటాం.
\end{defn}

$A \neq \emptyset$ అయితే, $A$పై ఏ అస్వావర్తన సంబంధమూ స్వావర్తన
సంబంధం కాదనీ, $A$పై ప్రతి ఏకదిశ సంబంధమూ ప్రతిసౌష్ఠవ సంబంధమేననీ
గమనించండి. అయితే స్వావర్తనంగానూ కాక అస్వావర్తనంగానూ కాని
$R \subseteq A^2$ సంబంధాలు ఉన్నాయి. అలాగే ఏకదిశగా కాని
ప్రతిసౌష్ఠవ సంబంధాలు కూడా ఉన్నాయి.

\end{document}
